\chapter{2010年辽宁卷（文）}

\section{单选题}

\begin{problem}
已知集合 \(U=\{1,3,5,7,9\}\)，\(A=\{1,5,7\}\)，则 \(\complement_U A=\)（\quad）

\choices
    {\(\{1,3\}\)}
    {\(\{3,7,9\}\)}
    {\(\{3,5,9\}\)}
    {\(\{3,9\}\)}
\end{problem}

\begin{answer}
D
\end{answer}

\begin{solution}
补集由全集中不属于 \(A\) 的元素组成，所以
\(\complement_U A=\{3,9\}\). 因此选 D.
\end{solution}

\begin{problem}
设 \(a\)，\(b\) 为实数，若复数 \(\frac{1+2\i}{a+b\i}=1+\i\)，则（\quad）

\choices
    {\(a=\frac{3}{2}\)，\(b=\frac{1}{2}\)}
    {\(a=3\)，\(b=1\)}
    {\(a=\frac{1}{2}\)，\(b=\frac{3}{2}\)}
    {\(a=1\)，\(b=3\)}
\end{problem}

\begin{answer}
A
\end{answer}

\begin{solution}
分母不为零，交叉相乘得
\(1+2\i=(1+\i)(a+b\i)=(a-b)+(a+b)\i\). 比较实部和虚部，得到
\(a-b=1\)，\(a+b=2\). 两式相加、相减，得 \(a=\frac{3}{2}\)，\(b=\frac{1}{2}\)，故选 A.
\end{solution}

\begin{problem}
设 \(S_n\) 为等比数列 \(\{a_n\}\) 的前 \(n\) 项和，已知 \(3S_3=a_4-2\)，\(3S_2=a_3-2\)，则公比 \(q=\)（\quad）

\choices
    {3}
    {4}
    {5}
    {6}
\end{problem}

\begin{answer}
B
\end{answer}

\begin{solution}
两式相减，并利用 \(S_3-S_2=a_3\)，得
\(3a_3=a_4-a_3\)，所以 \(a_4=4a_3\). 等比数列的公比按定义满足 \(q\ne0\). 若 \(a_3=0\)，由 \(a_3=a_1q^2\) 得 \(a_1=0\)，进而 \(S_2=0\)，这与 \(3S_2=a_3-2=-2\) 矛盾. 因而
\(q=\frac{a_4}{a_3}=4\)，选 B.
\end{solution}

\begin{problem}
已知 \(a>0\)，函数 \(f(x)=ax^2+bx+c\). 若 \(x_0\) 满足关于 \(x\) 的方程 \(2ax+b=0\)，则下列选项中的命题为假命题的是（\quad）

\choices
    {\(\exists x\in\R\)，\(f(x)\leqslant f(x_0)\)}
    {\(\exists x\in\R\)，\(f(x)\geqslant f(x_0)\)}
    {\(\forall x\in\R\)，\(f(x)\leqslant f(x_0)\)}
    {\(\forall x\in\R\)，\(f(x)\geqslant f(x_0)\)}
\end{problem}

\begin{answer}
C
\end{answer}

\begin{solution}
由 \(2ax_0+b=0\)，得 \(x_0=-\frac{b}{2a}\). 配方可得
\(f(x)=a\left(x+\frac{b}{2a}\right)^2+c-\frac{b^2}{4a}\)，\(f(x_0)=c-\frac{b^2}{4a}\).
因为 \(a>0\)，所以对任意 \(x\in\R\)，都有 \(f(x)\geqslant f(x_0)\)，且在 \(x=x_0\) 时取等号. 因此选项 A、B、D 都是真命题，而 \(\forall x\in\R\)，\(f(x)\leqslant f(x_0)\) 为假命题，选 C.
\end{solution}

\begin{problem}
如果执行下图所示的程序框图，输入 \(n=6\)，\(m=4\)，那么输出的 \(p\) 等于（\quad）

\bitmapfigure[max width=0.38\linewidth]{img/2010/liaoning_liberal/q5_fig1.png}

\choices
    {720}
    {360}
    {240}
    {120}
\end{problem}

\begin{answer}
B
\end{answer}

\begin{solution}
初始时 \(k=1\)，\(p=1\). 每次先计算 \(p\leftarrow p(n-m+k)\)，再判断 \(k<m\). 因此依次得到
\(p=1\times3=3\)，\(p=3\times4=12\)，\(p=12\times5=60\)，\(p=60\times6=360\).
最后一次计算后 \(k=4\)，不满足 \(k<4\)，程序输出 \(p=360\)，故选 B.
\end{solution}

\begin{problem}
设 \(\omega>0\)，函数 \(y=\sin\left(\omega x+\frac{\pi}{3}\right)+2\) 的图象向右平移 \(\frac{4\pi}{3}\) 个单位后与原图象重合，则 \(\omega\) 的最小值是（\quad）

\choices
    {\(\frac{2}{3}\)}
    {\(\frac{4}{3}\)}
    {\(\frac{3}{2}\)}
    {3}
\end{problem}

\begin{answer}
C
\end{answer}

\begin{solution}
函数的周期为 \(\frac{2\pi}{\omega}\). 平移距离 \(\frac{4\pi}{3}\) 后图象重合，说明该距离是周期的正整数倍，即存在正整数 \(k\)，使
\[
\frac{4\pi}{3}=k\frac{2\pi}{\omega}
\]
于是 \(\omega=\frac{3k}{2}\). 由于 \(k\geqslant1\)，\(\omega\) 的最小值为 \(\frac{3}{2}\)，选 C.
\end{solution}

\begin{problem}
设抛物线 \(y^2=8x\) 的焦点为 \(F\)，准线为 \(l\)，\(P\) 为抛物线上一点，\(PA\perp l\)，\(A\) 为垂足. 如果直线 \(AF\) 的斜率为 \(-\sqrt{3}\)，那么 \(|PF|=\)（\quad）

\choices
    {\(4\sqrt{3}\)}
    {8}
    {\(8\sqrt{3}\)}
    {16}
\end{problem}

\begin{answer}
B
\end{answer}

\begin{solution}
将 \(y^2=8x\) 写成 \(y^2=4\cdot2x\)，所以焦点为 \(F(2,0)\)，准线为 \(x=-2\). 设 \(P=(x,y)\)，则垂足为 \(A=(-2,y)\). 由直线 \(AF\) 的斜率，得
\[
\frac{0-y}{2-(-2)}=-\frac{y}{4}=-\sqrt{3}
\]
所以 \(y=4\sqrt{3}\). 代入抛物线方程得 \(x=\frac{y^2}{8}=6\). 因此
\[
|PF|=\sqrt{(6-2)^2+(4\sqrt{3})^2}=\sqrt{64}=8
\]
故选 B.
\end{solution}

\begin{problem}
平面上 \(O\)，\(A\)，\(B\) 三点不共线，设 \(\overrightarrow{OA}=\bs{a}\)，\(\overrightarrow{OB}=\bs{b}\)，则 \(\triangle OAB\) 的面积等于（\quad）

\choices
    {\(\sqrt{|\bs{a}|^2|\bs{b}|^2-(\bs{a}\cdot\bs{b})^2}\)}
    {\(\sqrt{|\bs{a}|^2|\bs{b}|^2+(\bs{a}\cdot\bs{b})^2}\)}
    {\(\frac{1}{2}\sqrt{|\bs{a}|^2|\bs{b}|^2-(\bs{a}\cdot\bs{b})^2}\)}
    {\(\frac{1}{2}\sqrt{|\bs{a}|^2|\bs{b}|^2+(\bs{a}\cdot\bs{b})^2}\)}
\end{problem}

\begin{answer}
C
\end{answer}

\begin{solution}
设 \(\bs{a}\) 与 \(\bs{b}\) 的夹角为 \(\theta\). 因为 \(O\)，\(A\)，\(B\) 不共线，所以 \(0<\theta<\pi\)，从而
\[
\left|\bs{a}\right|^2\left|\bs{b}\right|^2-(\bs{a}\cdot\bs{b})^2
=|\bs{a}|^2|\bs{b}|^2(1-\cos^2\theta)
=|\bs{a}|^2|\bs{b}|^2\sin^2\theta
\]
三角形面积为 \(\frac{1}{2}|\bs{a}||\bs{b}|\sin\theta\)，所以面积等于
\(\frac{1}{2}\sqrt{|\bs{a}|^2|\bs{b}|^2-(\bs{a}\cdot\bs{b})^2}\)，选 C.
\end{solution}

\begin{problem}
设双曲线的一个焦点为 \(F\)，虚轴的一个端点为 \(B\)，如果直线 \(FB\) 与该双曲线的一条渐近线垂直，那么此双曲线的离心率为（\quad）

\choices
    {\(\sqrt{2}\)}
    {\(\sqrt{3}\)}
    {\(\frac{\sqrt{3}+1}{2}\)}
    {\(\frac{\sqrt{5}+1}{2}\)}
\end{problem}

\begin{answer}
D
\end{answer}

\begin{solution}
不妨设双曲线的方程为 \(\frac{x^2}{a^2}-\frac{y^2}{b^2}=1\)，其中 \(a>0\)，\(b>0\). 记半焦距为 \(c\)，则 \(c^2=a^2+b^2\)，焦点和一个虚轴端点可取为 \(F(c,0)\)，\(B(0,b)\). 直线 \(FB\) 的斜率为 \(-\frac{b}{c}\)，渐近线的斜率可取 \(\frac{b}{a}\). 垂直条件给出
\(\left(-\frac{b}{c}\right)\left(\frac{b}{a}\right)=-1\)，\(b^2=ac\).
令离心率 \(e=\frac{c}{a}>1\)，由 \(b^2=c^2-a^2=ac\) 得
\(e^2-1=e\)，\(e^2-e-1=0\).
取大于 \(1\) 的根，得 \(e=\frac{1+\sqrt{5}}{2}\)，故选 D.
\end{solution}

\begin{problem}
设 \(2^a=5^b=m\)，且 \(\frac{1}{a}+\frac{1}{b}=2\)，则 \(m=\)（\quad）

\choices
    {\(\sqrt{10}\)}
    {10}
    {20}
    {100}
\end{problem}

\begin{answer}
A
\end{answer}

\begin{solution}
由 \(2^a=5^b=m\) 得 \(m>0\)，且 \(a=\log_2m\)，\(b=\log_5m\). 条件中的倒数存在，且其和为正，所以 \(m>1\). 利用换底公式，得
\[
\frac{1}{a}+\frac{1}{b}=\log_m2+\log_m5=\log_m10=2
\]
因此 \(m^2=10\). 结合 \(m>0\)，得 \(m=\sqrt{10}\)，选 A.
\end{solution}

\begin{problem}
已知 \(S\)，\(A\)，\(B\)，\(C\) 是球 \(O\) 表面上的点，\(SA\perp\) 平面 \(ABC\)，\(AB\perp BC\)，\(SA=AB=1\)，\(BC=\sqrt{2}\)，则球 \(O\) 的表面积等于（\quad）

\choices
    {\(4\pi\)}
    {\(3\pi\)}
    {\(2\pi\)}
    {\(\pi\)}
\end{problem}

\begin{answer}
A
\end{answer}

\begin{solution}
建立直角坐标系，令 \(A=(0,0,0)\)，平面 \(ABC\) 为 \(z=0\)，并取 \(B=(1,0,0)\). 由 \(AB\perp BC\) 和 \(BC=\sqrt{2}\)，可取 \(C=(1,\sqrt{2},0)\). 又因为 \(SA\perp\) 平面 \(ABC\) 且 \(SA=1\)，可取 \(S=(0,0,1)\).

设球心为 \(O=(u,v,w)\). 由 \(OA=OB\)，得
\(u^2+v^2+w^2=(u-1)^2+v^2+w^2\)，所以 \(u=\frac{1}{2}\). 再由 \(OA=OC\)，得
\(u^2+v^2+w^2=(u-1)^2+(v-\sqrt{2})^2+w^2\)，即 \(2u+2\sqrt{2}v=3\)，从而 \(v=\frac{1}{\sqrt{2}}\). 由 \(OA=OS\)，得
\(u^2+v^2+w^2=u^2+v^2+(w-1)^2\)，所以 \(w=\frac{1}{2}\).
所以
\[
R^2=OA^2=u^2+v^2+w^2=\frac14+\frac12+\frac14=1
\]
球的表面积为 \(4\pi R^2=4\pi\)，故选 A.
\end{solution}

\begin{problem}
已知点 \(P\) 在曲线 \(y=\frac{4}{\e^x+1}\) 上，\(\alpha\) 为曲线在点 \(P\) 处的切线的倾斜角，则 \(\alpha\) 的取值范围是（\quad）

\choices
    {\(\left[0,\frac{\pi}{4}\right)\)}
    {\(\left[\frac{\pi}{4},\frac{\pi}{2}\right)\)}
    {\(\left(\frac{\pi}{2},\frac{3\pi}{4}\right]\)}
    {\(\left[\frac{3\pi}{4},\pi\right)\)}
\end{problem}

\begin{answer}
D
\end{answer}

\begin{solution}
求导得
\[
y'=-\frac{4\e^x}{(\e^x+1)^2}
\]
令 \(t=\e^x>0\)，则 \(y'=-\frac{4t}{(t+1)^2}\). 因为 \((t+1)^2\geqslant4t\)，所以
\[
-1\leqslant y'<0
\]
当 \(t=1\)，即 \(x=0\) 时，\(y'=-1\)；对任意有限的 \(x\)，都有 \(t>0\)，故 \(y'\ne0\). 又当 \(t\to0^+\) 或 \(t\to+\infty\) 时，\(y'\to0^-\)，结合连续性可知 \(y'\) 的取值范围为 \([-1,0)\). 直线倾斜角取 \([0,\pi)\)，故由 \(\tan\alpha=y'\) 得
\(\alpha\in\left[\frac{3\pi}{4},\pi\right)\)，选 D.
\end{solution}

\section{填空题}

\begin{problem}
三张卡片上分别写上字母 \(E\)，\(E\)，\(B\)，将三张卡片随机地排成一行，恰好排成英文单词 \(BEE\) 的概率为\fillinblank{}.
\end{problem}

\begin{answer}
\(\frac{1}{3}\)
\end{answer}

\begin{solution}
若将两张写有 \(E\) 的卡片暂时区分为 \(E_1\)，\(E_2\)，三张卡片共有 \(3!=6\) 种等可能的排列. 排成字母序列 \(BEE\) 时，两张 \(E\) 的次序可以互换，恰有 \(2\) 种排列符合要求. 因此概率为 \(\frac{2}{6}=\frac{1}{3}\).
\end{solution}

\begin{problem}
设 \(S_n\) 为等差数列 \(\{a_n\}\) 的前 \(n\) 项和，若 \(S_3=3\)，\(S_6=24\)，则 \(a_9=\)\fillinblank{}.
\end{problem}

\begin{answer}
\(15\)
\end{answer}

\begin{solution}
设等差数列首项为 \(a_1\)，公差为 \(d\). 由求和公式，得
\(S_3=\frac{3}{2}(2a_1+2d)=3\)，\(S_6=\frac{6}{2}(2a_1+5d)=24\).
即 \(a_1+d=1\)，\(2a_1+5d=8\). 将第一式乘以 \(2\) 后相减，得 \(3d=6\)，所以 \(d=2\)，\(a_1=-1\). 因而
\(a_9=a_1+8d=-1+16=15\).
\end{solution}

\begin{problem}
已知 \(-1<x+y<4\) 且 \(2<x-y<3\)，则 \(z=2x-3y\) 的取值范围是\fillinblank{}. （答案用区间表示）
\end{problem}

\begin{answer}
\((3,8)\)
\end{answer}

\begin{solution}
令 \(u=x+y\)，\(v=x-y\)，则
\(x=\frac{u+v}{2}\)，\(y=\frac{u-v}{2}\)，\(z=2x-3y=\frac{-u+5v}{2}\).
由 \(-1<u<4\) 和 \(2<v<3\)，得
\(-4<-u<1\)，\(10<5v<15\). 因此
\[
3=\frac{-4+10}{2}<z<\frac{1+15}{2}=8
\]
两个开区间的线性和能取遍其中的所有值，所以 \(z\) 的取值范围为 \((3,8)\).
\end{solution}

\begin{problem}
如图，网格纸的小正方形的边长是 \(1\)，在其上用粗线画出了某多面体的三视图，则这个多面体最长的一条棱的长为\fillinblank{}.

\bitmapfigure[max width=0.56\linewidth]{img/2010/liaoning_liberal/q16_fig1.png}
\end{problem}

\begin{answer}
\(2\sqrt{3}\)
\end{answer}

\begin{solution}
把三幅投影按同一组坐标轴叠合，可还原为棱长为 \(2\) 的正方体切割后的一部分. 设三条相互垂直的投影方向为 \(x\)，\(y\)，\(z\)，并设三幅图中互相对应的斜边属于同一条棱，该棱两端点的坐标差分别为 \(\Delta x\)，\(\Delta y\)，\(\Delta z\). 每条斜边都是边长为 \(2\) 的网格正方形的对角线，因此三个正投影分别满足 \((\Delta x)^2+(\Delta y)^2=8\)，\((\Delta y)^2+(\Delta z)^2=8\)，\((\Delta z)^2+(\Delta x)^2=8\). 三式两两相减，得 \(|\Delta x|=|\Delta y|=|\Delta z|=2\). 由空间两点间的距离公式，这条棱的长度满足
\[
L^2=2^2+2^2+2^2=12
\]
所以 \(L=2\sqrt{3}\). 另一方面，三视图表明该多面体在三个坐标方向上的跨度均为 \(2\)，故任意一条棱的三个坐标分量绝对值都不超过 \(2\)，从而其长度不超过
\[
\sqrt{2^2+2^2+2^2}=2\sqrt{3}
\]
因此最长的一条棱的长为 \(2\sqrt{3}\).
\end{solution}

\section{解答题}

\begin{problem}
在 \(\triangle ABC\) 中，\(a\)，\(b\)，\(c\) 分别为内角 \(A\)，\(B\)，\(C\) 的对边，且 \(2a\sin A=(2b+c)\sin B+(2c+b)\sin C\).

\begin{enumerate}
\item 求 \(A\) 的大小；
\item 若 \(\sin B+\sin C=1\)，试判断 \(\triangle ABC\) 的形状.
\end{enumerate}
\end{problem}

\begin{solution}
\begin{enumerate}
\item 设三角形 \(ABC\) 的外接圆半径为 \(R\)，由正弦定理得 \(a=2R\sin A\)，\(b=2R\sin B\)，\(c=2R\sin C\). 将其代入已知等式，得 \(4R\sin^2 A=4R\left(\sin^2B+\sin^2C+\sin B\sin C\right)\)，从而 \(\sin^2 A=\sin^2B+\sin^2C+\sin B\sin C\). 又因为 \(A=\pi-(B+C)\)，所以
\[
\sin^2 A=\sin^2(B+C)=\sin^2B+\sin^2C-2\sin B\sin C\cos A
\]
因此 \(-2\sin B\sin C\cos A=\sin B\sin C\). 三角形内角 \(B\)，\(C\) 均在 \((0,\pi)\) 内，故 \(\sin B\sin C>0\)，于是 \(\cos A=-\frac{1}{2}\). 由 \(0<A<\pi\)，得 \(A=\frac{2\pi}{3}\).
\item 由上知 \(B+C=\frac{\pi}{3}\)，于是 \(\sin B+\sin C=2\sin\frac{B+C}{2}\cos\frac{B-C}{2}=\cos\frac{B-C}{2}\). 已知该和为 \(1\)，且 \(B\)，\(C\) 均为正数，所以只能有 \(B=C=\frac{\pi}{6}\). 因而 \(b=c\)，三角形 \(ABC\) 为等腰三角形，且 \(A=\frac{2\pi}{3}\)，\(B=C=\frac{\pi}{6}\).
\end{enumerate}
\end{solution}

\begin{problem}
为了比较注射 \(A\)，\(B\) 两种药物后产生的皮肤疱疹的面积，选 \(200\) 只家兔做试验，将这 \(200\) 只家兔随机地分成两组，每组 \(100\) 只，其中一组注射药物 \(A\)，另一组注射药物 \(B\). 下表 \(1\) 和表 \(2\) 分别是注射药物 \(A\) 和 \(B\) 后的试验结果. （疱疹面积单位：\(\mathrm{mm}^2\)）

\begin{center}
表 1：注射药物 \(A\) 后皮肤疱疹面积的频数分布表

\begin{tabular}{|c|c|c|c|c|}
\hline
疱疹面积 & \([60,65)\) & \([65,70)\) & \([70,75)\) & \([75,80)\) \\
\hline
频数 & 30 & 40 & 20 & 10 \\
\hline
\end{tabular}
\end{center}

\begin{center}
表 2：注射药物 \(B\) 后皮肤疱疹面积的频数分布表

\begin{tabular}{|c|c|c|c|c|c|}
\hline
疱疹面积 & \([60,65)\) & \([65,70)\) & \([70,75)\) & \([75,80)\) & \([80,85)\) \\
\hline
频数 & 10 & 25 & 20 & 30 & 15 \\
\hline
\end{tabular}
\end{center}

\begin{enumerate}
\item 完成下面频率分布直方图，并比较注射两种药物后疱疹面积的中位数大小；

\item 完成下面 \(2\times2\) 列联表，并回答能否有 \(99.9\%\) 的把握认为“注射药物 \(A\) 后的疱疹面积与注射药物 \(B\) 后的疱疹面积有差异”.
\end{enumerate}

\bitmapfigure[max width=0.92\linewidth]{img/2010/liaoning_liberal/q18_fig1.png}

\begin{center}
\begin{tabular}{|c|c|c|c|}
\hline
 & 疱疹面积小于 \(70\mathrm{~mm}^2\) & 疱疹面积不小于 \(70\mathrm{~mm}^2\) & 合计 \\
\hline
注射药物 \(A\) & \(a=\quad\) & \(b=\quad\) &  \\
\hline
注射药物 \(B\) & \(c=\quad\) & \(d=\quad\) &  \\
\hline
合计 &  &  & \(n=\quad\) \\
\hline
\end{tabular}
\end{center}

附：\(K^2=\dfrac{n(ad-bc)^2}{(a+b)(c+d)(a+c)(b+d)}\)

\begin{center}
\begin{tabular}{|c|c|c|c|c|c|}
\hline
\(P(K^2\geqslant k)\) & 0.100 & 0.050 & 0.025 & 0.010 & 0.001 \\
\hline
\(k\) & 2.706 & 3.841 & 5.024 & 6.635 & 10.828 \\
\hline
\end{tabular}
\end{center}
\end{problem}

\begin{solution}
\begin{enumerate}
\item 每组组距均为 \(5\)，每组样本数为 \(100\)，所以直方图的纵坐标 \(\text{频率}/\text{组距}\) 等于该组频数除以 \(500\). 药物 \(A\) 的四组矩形高依次为 \(\frac{30}{500}=0.06\)，\(\frac{40}{500}=0.08\)，\(\frac{20}{500}=0.04\)，\(\frac{10}{500}=0.02\)；药物 \(B\) 的五组矩形高依次为 \(\frac{10}{500}=0.02\)，\(\frac{25}{500}=0.05\)，\(\frac{20}{500}=0.04\)，\(\frac{30}{500}=0.06\)，\(\frac{15}{500}=0.03\). 因此按这些高度补全两幅直方图. 药物 \(A\) 的累计频数为 \(30\)，\(70\)，\(90\)，\(100\)，其中位数落在 \([65,70)\) 内；药物 \(B\) 的累计频数为 \(10\)，\(35\)，\(55\)，\(85\)，\(100\)，其中位数落在 \([70,75)\) 内，所以注射药物 \(B\) 后疱疹面积的中位数较大.
\item 由表中频数直接得到 \(a=30+40=70\)，\(b=20+10=30\)，\(c=10+25=35\)，\(d=20+30+15=65\). 因此列联表的行合计均为 \(100\)，列合计分别为 \(105\) 和 \(95\)，总计为 \(n=200\). 因而
\[
K^2=\frac{200(70\times65-30\times35)^2}{100\times100\times105\times95}\approx24.56
\]
由于 \(24.56>10.828\)，在显著性水平 \(0.001\) 下拒绝“注射药物 \(A\) 后与注射药物 \(B\) 后的疱疹面积没有差异”的假设，所以能有 \(99.9\%\) 的把握认为两种药物产生的疱疹面积有差异.
\end{enumerate}
\end{solution}

\begin{problem}
如图，棱柱 \(ABC-A_1B_1C_1\) 的侧面 \(BCC_1B_1\) 是菱形，\(B_1C\perp A_1B\).

\begin{enumerate}
\item 证明：平面 \(AB_1C\perp\) 平面 \(A_1BC_1\)；
\item 设 \(D\) 是 \(A_1C_1\) 上的点，且 \(A_1B\parallel\) 平面 \(B_1CD\)，求 \(A_1D:DC_1\) 的值.
\end{enumerate}

\bitmapfigure[max width=0.46\linewidth]{img/2010/liaoning_liberal/q19_fig1.png}
\end{problem}

\begin{solution}
\begin{enumerate}
\item 设 \(\overrightarrow{BA}=\bs{u}\)，\(\overrightarrow{BC}=\bs{v}\)，\(\overrightarrow{BB_1}=\bs{w}\). 取 \(B\) 为原点，则 \(A\)，\(C\)，\(B_1\)，\(A_1\)，\(C_1\) 的位置向量依次为 \(\bs{u}\)，\(\bs{v}\)，\(\bs{w}\)，\(\bs{u}+\bs{w}\)，\(\bs{v}+\bs{w}\). 因为四边形 \(BCC_1B_1\) 是菱形，所以 \(|\bs{v}|=|\bs{w}|\)，从而 \((\bs{v}-\bs{w})\cdot(\bs{v}+\bs{w})=|\bs{v}|^2-|\bs{w}|^2=0\)，即 \(B_1C\perp BC_1\). 由题设还有 \(B_1C\perp A_1B\). 设 \(E\) 为 \(BC_1\) 的中点，则菱形的两条对角线互相垂直且互相平分，所以 \(E\) 也在 \(B_1C\) 上. 过 \(E\) 作直线 \(g\parallel A_1B\)，则 \(g\) 在平面 \(A_1BC_1\) 内，且 \(B_1C\perp g\). 直线 \(BC_1\) 与 \(g\) 在 \(E\) 相交，因此 \(B_1C\) 垂直于平面 \(A_1BC_1\). 又直线 \(B_1C\) 在平面 \(AB_1C\) 内，故平面 \(AB_1C\perp\) 平面 \(A_1BC_1\).
\item 设 \(D=A_1+t(C_1-A_1)\)，其中 \(0<t<1\). 于是 \(D\) 的位置向量为 \((1-t)\bs{u}+t\bs{v}+\bs{w}\)，平面 \(B_1CD\) 内可取两条相交方向向量 \(\bs{v}-\bs{w}\) 与 \((1-t)\bs{u}+t\bs{v}\). 条件 \(A_1B\parallel\) 平面 \(B_1CD\) 等价于向量 \(\bs{u}+\bs{w}\) 可表示为这两条向量的线性组合，即存在实数 \(\lambda\)，\(\mu\)，使
\[
\bs{u}+\bs{w}=\lambda(\bs{v}-\bs{w})+\mu\bigl((1-t)\bs{u}+t\bs{v}\bigr)
\]
比较 \(\bs{u}\)，\(\bs{v}\)，\(\bs{w}\) 的系数，得到 \(\mu(1-t)=1\)，\(\lambda+\mu t=0\)，\(-\lambda=1\). 所以 \(\lambda=-1\)，\(\mu t=1\)，且 \(\mu(1-t)=1\)，从而 \(t=\frac{1}{2}\). 因此 \(A_1D:DC_1=t:(1-t)=1:1\).
\end{enumerate}
\end{solution}

\begin{problem}
设 \(F_1\)，\(F_2\) 分别为椭圆 \(C:\frac{x^2}{a^2}+\frac{y^2}{b^2}=1\)（\(a>b>0\)）的左，右焦点，过 \(F_2\) 的直线 \(l\) 与椭圆 \(C\) 相交于 \(A\)，\(B\) 两点，直线 \(l\) 的倾斜角为 \(60^\circ\)，\(F_1\) 到直线 \(l\) 的距离为 \(2\sqrt{3}\).

\begin{enumerate}
\item 求椭圆 \(C\) 的焦距；
\item 如果 \(\overrightarrow{AF_2}=2\overrightarrow{F_2B}\)，求椭圆 \(C\) 的方程.
\end{enumerate}
\end{problem}

\begin{solution}
\begin{enumerate}
\item 设椭圆的半焦距为 \(c\)，则 \(F_1=(-c,0)\)，\(F_2=(c,0)\). 直线 \(l\) 过 \(F_2\)，倾斜角为 \(60^\circ\)，其方程为 \(y=\sqrt{3}(x-c)\)，即 \(\sqrt{3}x-y-\sqrt{3}c=0\). 因此 \(F_1\) 到直线 \(l\) 的距离为 \(\frac{2\sqrt{3}c}{2}=\sqrt{3}c\). 由题意 \(\sqrt{3}c=2\sqrt{3}\)，得 \(c=2\)，所以椭圆的焦距 \(F_1F_2=2c=4\).
\item 由 \(\overrightarrow{AF_2}=2\overrightarrow{F_2B}\)，取直线 \(l\) 的单位方向向量 \(\bs{e}=\left(\frac{1}{2},\frac{\sqrt{3}}{2}\right)\)，并设 \(\overrightarrow{F_2B}=t\bs{e}\)，其中 \(t>0\). 则 \(\overrightarrow{AF_2}=2t\bs{e}\)，故
\(A=\left(2-t,-\sqrt{3}t\right)\)，\(B=\left(2+\frac{t}{2},\frac{\sqrt{3}t}{2}\right)\).
令直线上的点为 \(X(s)=F_2+s\bs{e}\)，将其代入椭圆方程，得到
\[
\frac{(2+s/2)^2}{a^2}+\frac{3s^2}{4b^2}=1
\]
即
\[
\left(\frac{1}{4a^2}+\frac{3}{4b^2}\right)s^2+\frac{2}{a^2}s+\frac{4}{a^2}-1=0
\]
该方程的两个根分别为 \(-2t\) 与 \(t\). 记 \(D=b^2+3a^2\)，由韦达定理，根的和给出 \(-t=-\frac{8b^2}{D}\)，即 \(t=\frac{8b^2}{D}\)；根的积给出 \(-2t^2=\frac{4b^2(4-a^2)}{D}\). 又因 \(c^2=a^2-b^2=4\)，所以 \(b^2=a^2-4>0\). 将 \(t=\frac{8b^2}{D}\) 代入根的积关系，得 \(-\frac{128b^4}{D^2}=-\frac{4b^4}{D}\)，从而 \(D=32\). 因此 \(b^2+3a^2=32\)，结合 \(b^2=a^2-4\)，得 \(a^2=9\)，\(b^2=5\). 所以椭圆 \(C\) 的方程为 \(\dfrac{x^2}{9}+\dfrac{y^2}{5}=1\).
\end{enumerate}
\end{solution}

\begin{problem}
已知函数 \(f(x)=(a+1)\ln x+ax^2+1\).

\begin{enumerate}
\item 讨论函数 \(f(x)\) 的单调性；
\item 设 \(a\leqslant -2\)，证明：对任意 \(x_1\)，\(x_2\in(0,+\infty)\)，\(|f(x_1)-f(x_2)|\geqslant 4|x_1-x_2|\).
\end{enumerate}
\end{problem}

\begin{solution}
\begin{enumerate}
\item 函数的定义域为 \((0,+\infty)\)，且
\[
f'(x)=\frac{a+1}{x}+2ax=\frac{a+1+2ax^2}{x}
\]
由于 \(x>0\)，\(f'(x)\) 的符号由 \(a+1+2ax^2\) 决定. 当 \(a\geqslant0\) 时，\(a+1+2ax^2>0\)，所以 \(f(x)\) 在 \((0,+\infty)\) 上单调增加；当 \(a\leqslant-1\) 时，\(a+1+2ax^2<0\)，所以 \(f(x)\) 在 \((0,+\infty)\) 上单调减少；当 \(-1<a<0\) 时，方程 \(a+1+2ax^2=0\) 的正根为 \(x_0=\sqrt{-\dfrac{a+1}{2a}}\)，且 \(f'(x)>0\) 在 \(\left(0,x_0\right)\) 上成立，\(f'(x)<0\) 在 \(\left(x_0,+\infty\right)\) 上成立，所以 \(f(x)\) 在 \(\left(0,\sqrt{-\dfrac{a+1}{2a}}\right)\) 上单调增加，在 \(\left(\sqrt{-\dfrac{a+1}{2a}},+\infty\right)\) 上单调减少.
\item 设 \(u=-a\)，则 \(u\geqslant2\)，且
\[
f'(x)=-\frac{u-1}{x}-2ux
\]
由基本不等式，\(-f'(x)=\frac{u-1}{x}+2ux\geqslant2\sqrt{2u(u-1)}\geqslant4\)，所以 \(f'(x)\leqslant-4\) 对任意 \(x>0\) 成立. 当 \(x_1=x_2\) 时，不等式显然成立；当 \(x_1\ne x_2\) 时，由拉格朗日中值定理，存在 \(\xi\) 严格介于 \(x_1\)，\(x_2\) 之间，使 \(f(x_1)-f(x_2)=f'(\xi)(x_1-x_2)\). 因而 \(|f(x_1)-f(x_2)|=|f'(\xi)|\,|x_1-x_2|\geqslant4|x_1-x_2|\).
\end{enumerate}
\end{solution}

\section{选考题：解答题}

\begin{problem}
如图，\(\triangle ABC\) 的角平分线 \(AD\) 的延长线交它的外接圆于点 \(E\).

\begin{enumerate}
\item 证明：\(\triangle ABE\sim\triangle ADC\)；
\item 若 \(\triangle ABC\) 的面积 \(S=\frac{1}{2}AD\cdot AE\)，求 \(\angle BAC\) 的大小.
\end{enumerate}

\bitmapfigure[max width=0.34\linewidth]{img/2010/liaoning_liberal/q22_fig1.png}
\end{problem}

\begin{solution}
\begin{enumerate}
\item 记 \(\angle BAC=A\)，\(\angle ABC=B\)，\(\angle ACB=C\). 因为 \(AD\) 是角平分线且 \(A\)，\(D\)，\(E\) 共线，所以 \(\angle BAE=\angle DAC=\frac{A}{2}\). 又因为 \(A\)，\(B\)，\(C\)，\(E\) 共圆，且 \(B\)，\(C\) 在弦 \(AE\) 的两侧，有 \(\angle ABE=\pi-\angle ACE\). 在三角形 \(ACE\) 中，\(\angle CAE=\frac{A}{2}\)，并由圆周角定理 \(\angle AEC=\angle ABC=B\)，所以 \(\angle ACE=\pi-\frac{A}{2}-B\)，从而 \(\angle ABE=\frac{A}{2}+B\). 另一方面，\(D\) 在 \(BC\) 上，故 \(\angle ACD=C\)，于是 \(\angle ADC=\pi-\frac{A}{2}-C=\frac{A}{2}+B\). 因此 \(\angle BAE=\angle DAC\)，\(\angle ABE=\angle ADC\)，所以 \(\triangle ABE\sim\triangle ADC\).
\item 由相似三角形的对应边关系，得 \(\frac{AB}{AD}=\frac{AE}{AC}\)，即 \(AD\cdot AE=AB\cdot AC\). 另一方面，\(S=\frac{1}{2}AB\cdot AC\sin A\)，题设又给出 \(S=\frac{1}{2}AD\cdot AE=\frac{1}{2}AB\cdot AC\). 因为 \(AB\cdot AC>0\)，所以 \(\sin A=1\). 由 \(0<A<\pi\)，得 \(\angle BAC=A=\frac{\pi}{2}=90^\circ\).
\end{enumerate}
\end{solution}

\begin{problem}
已知 \(P\) 为半圆 \(C:\begin{cases}
x=\cos\theta,\\
y=\sin\theta,
\end{cases}\) （\(\theta\) 为参数，\(0\leqslant\theta\leqslant\pi\)）上的点，点 \(A\) 的坐标为 \((1,0)\)，\(O\) 为坐标原点，点 \(M\) 在射线 \(OP\) 上，线段 \(OM\) 与 \(C\) 的弧 \(\widehat{AP}\) 的长度均为 \(\frac{\pi}{3}\).

\begin{enumerate}
\item 以 \(O\) 为极点，\(x\) 轴的正半轴为极轴建立极坐标系，求点 \(M\) 的极坐标；
\item 求直线 \(AM\) 的参数方程.
\end{enumerate}
\end{problem}

\begin{solution}
\begin{enumerate}
\item 半圆 \(C\) 的半径为 \(1\)，点 \(A\) 对应参数 \(\theta=0\). 从 \(A\) 到 \(P\) 的弧长等于参数角 \(\theta\)，由题意 \(\theta=\frac{\pi}{3}\)，所以 \(P=\left(\cos\frac{\pi}{3},\sin\frac{\pi}{3}\right)=\left(\frac{1}{2},\frac{\sqrt{3}}{2}\right)\). 又因 \(M\) 在射线 \(OP\) 上且 \(OM=\frac{\pi}{3}\)，故点 \(M\) 的极坐标为 \(\left(\frac{\pi}{3},\frac{\pi}{3}\right)\).
\item 将点 \(M\) 化为直角坐标，得 \(M=\left(\frac{\pi}{3}\cos\frac{\pi}{3},\frac{\pi}{3}\sin\frac{\pi}{3}\right)=\left(\frac{\pi}{6},\frac{\sqrt{3}\pi}{6}\right)\). 于是 \(\overrightarrow{AM}=\left(\frac{\pi}{6}-1,\frac{\sqrt{3}\pi}{6}\right)\)，直线 \(AM\) 的参数方程为
\[
\begin{cases}
x=1+\left(\dfrac{\pi}{6}-1\right)t,\\
y=\dfrac{\sqrt{3}\pi}{6}t,
\end{cases}
\quad(t\in\mathbb{R})
\]
\end{enumerate}
\end{solution}

\begin{problem}
已知 \(a\)，\(b\)，\(c\) 均为正数，证明：\(a^2+b^2+c^2+\left(\frac{1}{a}+\frac{1}{b}+\frac{1}{c}\right)^2\geqslant 6\sqrt{3}\)，并确定 \(a\)，\(b\)，\(c\) 为何值时，等号成立.
\end{problem}

\begin{solution}
设 \(p=ab+bc+ca\)，\(q=\frac{1}{ab}+\frac{1}{bc}+\frac{1}{ca}\). 由 \(a^2+b^2+c^2\geqslant ab+bc+ca=p\)，以及 \((x+y+z)^2\geqslant3(xy+yz+zx)\)，取 \(x=\frac{1}{a}\)，\(y=\frac{1}{b}\)，\(z=\frac{1}{c}\)，得
\[
\left(\frac{1}{a}+\frac{1}{b}+\frac{1}{c}\right)^2\geqslant3q
\]
又由柯西不等式，\(pq=(ab+bc+ca)\left(\frac{1}{ab}+\frac{1}{bc}+\frac{1}{ca}\right)\geqslant9\). 因此
\[
a^2+b^2+c^2+\left(\frac{1}{a}+\frac{1}{b}+\frac{1}{c}\right)^2\geqslant p+3q\geqslant2\sqrt{3pq}\geqslant6\sqrt{3}
\]
要使等号成立，首先需有 \(a=b=c\). 设 \(a=b=c=t>0\)，则 \(p=3t^2\)，\(q=\frac{3}{t^2}\)，而 \(p+3q\geqslant2\sqrt{3pq}\) 取等号的条件为 \(p=3q\)，故 \(3t^2=\frac{9}{t^2}\)，即 \(t^4=3\). 所以当 \(a=b=c=\sqrt[4]{3}\) 时等号成立.
\end{solution}
